\documentclass[12pt]{article}
\usepackage{amsmath, amssymb}
\usepackage[margin=1in]{geometry}
\setlength{\parskip}{6pt}
\title{QUBO Formulation: Network Intrusion Feature Selection}
\date{}
\begin{document}
\maketitle

\subsection*{Network Intrusion Feature Selection}

\paragraph{Problem Statement.}
Given a set of $N$ candidate network-traffic features, each associated with a detection relevance score $s_i$ and pairwise redundancy penalties $r_{ij}$, the goal is to select exactly $K$ features that maximize the total detection relevance while minimizing the total redundancy among the selected features.

\paragraph{Decision Variables.}
For each candidate feature $i \in \{0, 1, \dots, N-1\}$, define a binary decision variable $x_i \in \{0,1\}$, where $x_i = 1$ if feature $i$ is selected for the intrusion detection model, and $x_i = 0$ otherwise.

\paragraph{QUBO Derivation.}
Step 1 --- The original maximization objective is $\max_{x} \sum_{i=0}^{N-1} s_i x_i - \sum_{0 \le i < j \le N-1} r_{ij} x_i x_j$. To cast this as a minimization problem suitable for QUBO, we negate the objective:
\begin{align}
  \min_{x} \left( -\sum_{i=0}^{N-1} s_i x_i + \sum_{0 \le i < j \le N-1} r_{ij} x_i x_j \right).
\end{align}

Step 2 --- The selection constraint requires exactly $K$ features to be chosen:
\begin{align}
  \sum_{i=0}^{N-1} x_i = K.
\end{align}

Step 3 --- We enforce this constraint via a quadratic penalty term with weight $\lambda > 0$:
\begin{align}
  P(x) = \lambda \left( \sum_{i=0}^{N-1} x_i - K \right)^2.
\end{align}
Expanding the square and applying the idempotent property $x_i^2 = x_i$ for binary variables yields:
\begin{align}
  P(x) &= \lambda \left( \sum_{i=0}^{N-1} x_i^2 - 2K \sum_{i=0}^{N-1} x_i + K^2 + \sum_{i \neq j} x_i x_j \right) \\
  &= \lambda \left( \sum_{i=0}^{N-1} x_i - 2K \sum_{i=0}^{N-1} x_i + K^2 + \sum_{i \neq j} x_i x_j \right) \\
  &= \lambda (1 - 2K) \sum_{i=0}^{N-1} x_i + \lambda K^2 + \lambda \sum_{i \neq j} x_i x_j.
\end{align}
Noting that $\sum_{i \neq j} x_i x_j = 2 \sum_{0 \le i < j \le N-1} x_i x_j$, the penalty simplifies to:
\begin{align}
  P(x) = \lambda (1 - 2K) \sum_{i=0}^{N-1} x_i + 2\lambda \sum_{0 \le i < j \le N-1} x_i x_j + \lambda K^2.
\end{align}

Step 4 --- Combining the negated objective and the penalty term gives the full QUBO Hamiltonian $H(x)$:
\begin{align}
  H(x) &= \left( -\sum_{i=0}^{N-1} s_i x_i + \sum_{0 \le i < j \le N-1} r_{ij} x_i x_j \right) + \left( \lambda (1 - 2K) \sum_{i=0}^{N-1} x_i + 2\lambda \sum_{0 \le i < j \le N-1} x_i x_j + \lambda K^2 \right) \\
  &= \sum_{i=0}^{N-1} \left( \lambda (1 - 2K) - s_i \right) x_i + \sum_{0 \le i < j \le N-1} \left( r_{ij} + 2\lambda \right) x_i x_j + \lambda K^2.
\end{align}

Step 5 --- Reading off the coefficients from the expanded Hamiltonian, the upper-triangular entries of the QUBO matrix $Q$ and the constant offset $c$ are given by:
\begin{align}
  Q_{i,i} &= \lambda (1 - 2K) - s_i, & \text{for } i \in \{0, \dots, N-1\}, \\
  Q_{i,j} &= r_{ij} + 2\lambda, & \text{for } 0 \le i < j \le N-1, \\
  c &= \lambda K^2.
\end{align}

\paragraph{Penalty Weight Justification.}
The penalty weight $\lambda$ must be chosen sufficiently large to ensure that any violation of the cardinality constraint incurs an energy cost strictly greater than the maximum possible improvement in the original objective function. The maximum possible gain in the objective from arbitrarily changing the selection of features is bounded by the sum of the absolute values of all linear and quadratic coefficients. Specifically, for any feature $i$, the maximum objective gain achievable by violating the constraint is at most $\max_i |s_i| + \sum_{j \neq i} |r_{ij}|$. Therefore, choosing $\lambda$ such that
\begin{align}
  \lambda > \max_{i \in \{0,\dots,N-1\}} |s_i| + \sum_{j \neq i} |r_{ij}|
\end{align}
guarantees that the penalty term dominates any feasible objective improvement, strictly enforcing the constraint at optimality.

\paragraph{Correctness Argument.}
Minimizing $H(x)$ over $\{0,1\}^N$ recovers the optimal feature selection because (i) all feasible solutions satisfying $\sum_i x_i = K$ incur zero penalty, reducing $H(x)$ exactly to the negated original objective; and (ii) any infeasible solution violates the constraint by at least one unit, incurring a penalty of at least $\lambda$. By the choice of $\lambda$ in the justification section, the energy of any infeasible configuration strictly exceeds the energy of the optimal feasible configuration. Consequently, the global minimum of $H(x)$ must correspond to a feasible assignment that minimizes the negated objective, which is equivalent to maximizing the original objective.

\paragraph{Illustrative Example.}
Consider a concrete instance with $N=3$ candidate features and $K=2$ required selections. Let the relevance scores be $s_0, s_1, s_2$ and redundancy penalties be $r_{01}, r_{02}, r_{12}$. Substituting $N=3$ and $K=2$ into the general formulas yields the explicit Hamiltonian:
\begin{align}
  H(x) &= \sum_{i=0}^{2} \left( -3\lambda - s_i \right) x_i + \left( r_{01} + 2\lambda \right) x_0 x_1 + \left( r_{02} + 2\lambda \right) x_0 x_2 + \left( r_{12} + 2\lambda \right) x_1 x_2 + 4\lambda.
\end{align}
The corresponding QUBO matrix entries and offset are:
\begin{align}
  Q_{0,0} &= -3\lambda - s_0, & Q_{0,1} &= r_{01} + 2\lambda, \\
  Q_{1,1} &= -3\lambda - s_1, & Q_{1,2} &= r_{12} + 2\lambda, \\
  Q_{2,2} &= -3\lambda - s_2, & Q_{0,2} &= r_{02} + 2\lambda, \\
  c &= 4\lambda.
\end{align}
For a specific numerical instantiation, let $s_0=5, s_1=3, s_2=4$, $r_{01}=1, r_{02}=2, r_{12}=1$, and $\lambda=100$. The Hamiltonian evaluates to:
\begin{align}
  H(x) &= -105 x_0 - 103 x_1 - 104 x_2 + 201 x_0 x_1 + 202 x_0 x_2 + 201 x_1 x_2 + 400.
\end{align}
Enumerating all $\binom{3}{2}=3$ feasible assignments ($x \in \{(1,1,0), (1,0,1), (0,1,1)\}$) and comparing against infeasible ones confirms that the minimum energy occurs at $x^* = (1,1,0)$ with $H(x^*) = 195$, corresponding to selecting features 0 and 1 for a maximum original objective score of $5+3-1=7$.

\end{document}